\providecommand{\vf}[1]{\textcolor{blue}{#1}}
\chapter{Future Work}\label{ch:future}

The attack of Chapter~\ref{ch:attack} maps one hidden state at a time, taken from one block of the target model.
Future work can extend it in two ways: by reading hidden states from several blocks, and by reading the hidden states of a prompt as a sequence.
A third direction studies which parts of a hidden state carry the prompt, and what that means for defenses.

\paragraph{Hidden states from several blocks.} Each block represents the same token in a different way, together with its context in the prompt.
The first question is whether a map that reads several blocks together recovers tokens that a single block loses.
The test compares reconstruction from each block alone with reconstruction from the blocks combined, at the same leak of \vf{90} rows and with the hub frozen, so that the number of blocks is the one difference.
Reading several blocks also changes what the adversary observes.
A server that runs the remaining blocks computes the hidden states of deeper blocks itself, whereas the adversary of Chapter~\ref{ch:attack} holds no weights of the target model.
That work therefore has to state which blocks the adversary observes, and it has to keep the question about information across blocks apart from the attack on one transmitted block.

\paragraph{The sequence of tokens.} The hub and the map send each hidden state into the shared space independently, and only the frozen decoder reads the sequence.
Yet each hidden state is computed from the tokens before it, so neighboring hidden states carry overlapping information about one token.
A map that reads a window of hidden states could draw on that overlap where a single hidden state leaves a token ambiguous.
The examples of Chapter~\ref{ch:attack} show wrong names, dates, and numbers inside otherwise correct reconstructions, and those tokens are the natural target.
That information is already present in the transmitted hidden states, and the current attack does not exploit it.

\paragraph{Representations and privacy.} The hub shows that one shared space holds the prompt information of several models, and it does not show where in a hidden state that information sits.
The open question is which components of a hidden state carry the identity of a token, and which carry only what the server's blocks need to continue the computation.
Candidate components can be removed or perturbed before reconstruction, and the change in word-F1 then separates the components that reconstruction depends on from those that merely correlate with the prompt.
Defenses that perturb token embeddings on the client trade reconstruction against the quality of the model's output~\cite{mai2024split,nguyen2026noir}.
A defense that knows where the content of a prompt sits could suppress those components and keep the ones that inference needs.
Any such defense has to face an adversary that fits its map again on the defended hidden states, with the same leak.
Without that second fit, a lower score shows only that the old map no longer reads the hidden states, not that the information is gone.
The same study extends the evidence to further target models and blocks, since the present results cover one undefended target model.

Each direction asks the same question from another side: how much of a prompt its hidden states carry, and how much of that a client can withhold.
